\documentclass[11pt]{article}
\usepackage[T1]{fontenc}
\usepackage[utf8]{inputenc}
\usepackage{lmodern}
\usepackage[margin=1in]{geometry}
\usepackage{amsmath,amssymb,amsthm}
\usepackage{microtype}
\usepackage[colorlinks=true,linkcolor=blue,citecolor=blue,urlcolor=blue]{hyperref}
\hypersetup{
  pdftitle={Constructive Grounded Hyperset Theory: Anti-Foundation, Bisimulation Quotients, and SCC Condensation for Hypercomputational Oracles},
  pdfauthor={Tyler Roost},
  pdfsubject={Non-well-founded sets, Aczel AFA, Paige-Tarjan bisimulation, SCC condensation, and superintelligence oracles},
  pdfkeywords={Hyperset theory, Anti-Foundation Axiom, Accessible Pointed Graphs, bisimulation, SCC condensation, Tarjan, real analysis bridge, superintelligence, Verifier Standard}
}
\usepackage{enumitem}
\setlist{nosep}
\newtheorem{theorem}{Theorem}[section]
\newtheorem{proposition}[theorem]{Proposition}
\newtheorem{lemma}[theorem]{Lemma}
\theoremstyle{definition}
\newtheorem{definition}[theorem]{Definition}
\newtheorem{example}[theorem]{Example}
\newtheorem{remark}[theorem]{Remark}

\newcommand{\N}{\mathbb N}
\newcommand{\Q}{\mathbb Q}
\newcommand{\R}{\mathbb R}
\newcommand{\APG}{\operatorname{APG}}
\newcommand{\code}[1]{\ulcorner #1\urcorner}
\newcommand{\Fix}{\operatorname{Fix}}
\newcommand{\SCC}{\operatorname{SCC}}

\title{Constructive Grounded Hyperset Theory:\\Anti-Foundation, Bisimulation Quotients, and SCC Condensation\\for Hypercomputational Oracles}
\author{Tyler Roost}
\date{Preprint, version 0.2.0 --- September 18, 2026}

\begin{document}
\maketitle

\begin{abstract}
We present a constructive, computational formulation of non-well-founded set theory grounded in Aczel's Anti-Foundation Axiom (AFA), explicit Accessible Pointed Graph (APG) representations, Paige-Tarjan relational partition refinement, and Tarjan Strongly Connected Component (SCC) condensation. The core impetus of this framework is to discover routes toward definably complete representations capable of handling hypercomputation-like oracle classes expected of superintelligences. By replacing ill-founded membership trees with bisimulation-invariant coalgebraic state machines, the theory provides a rigorous, paradox-free setting for circular structures, recursive self-models, and transfinite dynamical attractors. We introduce an explicit dual-abstraction paradigm distinguishing bisimulation minimization (which preserves observational behavior) from SCC condensation (which contracts cyclic components into acyclic Directed Acyclic Graphs, deciding trajectory confluence via terminal SCC analysis). Furthermore, we construct an explicit bridge from hyperset quotient spaces to computable real analysis, metric completions, and Verifier Standard (VSTD) graph assurance dynamics. The four theorems below carry prose proofs resting on established algorithmic results; no theorem here is machine-checked, and VSTD receipt emission is described as the intended verification interface rather than implemented, so no result in this paper rests on receipt evidence.
\end{abstract}

\noindent\textbf{Keywords:} Hypersets; Anti-Foundation Axiom; Accessible Pointed Graphs; bisimulation quotients; Paige-Tarjan; Strongly Connected Component condensation; confluence; real analysis bridge; superintelligence oracles; Verifier Standard.

\section{Introduction and Foundational Impetus}
Standard set-theoretic foundations, such as Zermelo--Fraenkel set theory with the Axiom of Choice (abbreviated ZFC), enforce the Axiom of Foundation (Regularity), stipulating that every non-empty set $x$ contains an $\in$-minimal element. While Regularity precludes circular membership chains such as $x \in x$ or $x \in y \in x$, it severely complicates the formal representation of reflexive and self-referential systems. In computational cognitive architectures and models of hypothetical superintelligences, reflexive self-containment is an operational necessity: an autonomous agent must maintain a model of its own cognitive state, its execution environment, and its predictive cycles without generating an infinite regress of ungrounded meta-representations.

Furthermore, advanced cognitive agents are expected to operate relative to hypercomputational oracle classes (such as the halting oracle $\mathbf{0}'$, hyperarithmetic sets $\Delta^1_1$, and transfinite ordinal jumps). A sound foundation must therefore represent both circular self-models and transfinite oracle ladders without collapsing into semantic inconsistency or unearned epistemic authority.

\emph{Grounded Hyperset Theory} provides this foundation by synthesizing:
\begin{enumerate}
\item \textbf{Aczel's Anti-Foundation Axiom (AFA):} Stating that every Accessible Pointed Graph (APG) has a unique decorative solution as a set \cite{aczel}.
\item \textbf{Bisimulation Quotients:} Computing canonical bisimulation classes in $O(|E| \log |V|)$ time via Paige-Tarjan partition refinement \cite{paige_tarjan}.
\item \textbf{Strongly Connected Component (SCC) Condensation:} Contracting cyclic state subgraphs via Tarjan's linear-time algorithm into an acyclic Directed Acyclic Graph (DAG), enabling exact confluence decisions via terminal SCC counting.
\item \textbf{The Real Analysis Bridge:} Connecting non-well-founded set representations to metric completions, continuous intervals, and dynamical flows over $\R$.
\item \textbf{Stratified Oracle Calibration:} Grounding each level of hyperset abstraction within the Verifier Standard (VSTD) claim discipline.
\end{enumerate}

\section{Accessible Pointed Graphs and Aczel's AFA}
In constructive hyperset theory, sets are represented by directed graphs whose nodes correspond to sets and whose directed edges represent membership.

\begin{definition}[Accessible Pointed Graph (APG)]
An \emph{Accessible Pointed Graph} (APG) is a triple $G = (V, E, r)$ where $V$ is a finite set of vertices, $E \subseteq V \times V$ is a set of directed edges, and $r \in V$ is the root vertex such that for every $v \in V$, there exists a directed path from $r$ to $v$. A node $v$ is a child of $u$ if $(u, v) \in E$.
\end{definition}

\begin{definition}[Decoration of an APG]
A \emph{decoration} of an APG $G = (V, E, r)$ is a function $d$ assigning to each node $v \in V$ a set such that
\[
d(v) = \{d(u) : (v, u) \in E\}.
\]
Aczel's Anti-Foundation Axiom (AFA) asserts that every APG has a unique decoration. The set \emph{depicted} by $G$ is $d(r)$.
\end{definition}

\begin{example}[The Quine Atom $\Omega$]
Consider the single-node self-loop graph $G_\Omega = (\{r\}, \{(r, r)\}, r)$. Its unique decoration satisfies $d(r) = \{d(r)\}$. Thus $\Omega = \{\Omega\}$ is a well-defined, unique hyperset.
\end{example}

\section{Computable Bisimulation via Paige-Tarjan Refinement}
Two APGs depict the same hyperset under AFA if and only if they are bisimilar.

\begin{definition}[Bisimulation Relation]
Given graphs $G_1 = (V_1, E_1, r_1)$ and $G_2 = (V_2, E_2, r_2)$, a relation $R \subseteq V_1 \times V_2$ is a \emph{bisimulation} if $r_1 R r_2$ and whenever $u R v$:
\begin{enumerate}
\item $\forall u' \in E_1(u),\, \exists v' \in E_2(v)$ such that $u' R v'$.
\item $\forall v' \in E_2(v),\, \exists u' \in E_1(u)$ such that $u' R v'$.
\end{enumerate}
\end{definition}

\begin{theorem}[Constructive Quotient and Canonical Minimization]
Let $G = (V, E, r)$ be a finite APG. The maximal bisimulation equivalence relation $\approx$ on $V$ can be computed constructively in time $O(|E| \log |V|)$ using Paige-Tarjan relational partition refinement. The quotient graph $G^* = G / \approx$ is the unique minimal strongly extensional representation of the hyperset depicted by $G$.
\end{theorem}
\begin{proof}
Paige-Tarjan partition refinement iteratively splits compound blocks in a partition $P$ of $V$ with respect to splitters until stability is achieved. Since $V$ is finite, the process terminates deterministically and produces the coarsest relational partition respecting edge transitions, which coincides exactly with maximal bisimulation \cite{paige_tarjan}.
\end{proof}

\section{Strongly Connected Component (SCC) Condensation and Confluence}
While Paige-Tarjan bisimulation quotienting yields the minimal behavioral representation, it preserves cyclic edges whenever non-trivial transitions occur. However, macroscopic structural analysis, stratified stage assignment, and confluence testing require abstracting away cyclic recurrence to expose the underlying dependency structure.

\subsection{Tarjan SCC Condensation}
\begin{definition}[SCC Quotient Graph]
Let $G = (V, E, r)$ be an APG. A \emph{Strongly Connected Component} (SCC) is a maximal subset $C \subseteq V$ such that every pair of vertices in $C$ is mutually reachable along directed paths. The \emph{SCC condensation} is the quotient graph
\[
G_{\SCC} = (V_{\SCC}, E_{\SCC}, [r]),
\]
where $V_{\SCC} = \{[v] : v \in V\}$ is the set of equivalence classes under mutual reachability, and $([u], [v]) \in E_{\SCC}$ if and only if $[u] \ne [v]$ and there exist $u' \in [u], v' \in [v]$ with $(u', v') \in E$.
\end{definition}

\begin{theorem}[Acyclicity of the SCC Condensation]
For any APG $G$, the condensation graph $G_{\SCC}$ is a strictly acyclic Directed Acyclic Graph (DAG).
\end{theorem}
\begin{proof}
Suppose $G_{\SCC}$ contained a directed cycle $[v_1] \to [v_2] \to \dots \to [v_k] \to [v_1]$ with $k \ge 2$. Then every node in $[v_1]$ can reach every node in $[v_k]$, and vice versa. By maximality of SCCs, $\bigcup_{i=1}^k [v_i]$ forms a single strongly connected component, so $[v_1] = [v_2] = \dots = [v_k]$, contradicting $[u] \ne [v]$.
\end{proof}

In our software realization (\texttt{abstraction.py}), Tarjan's linear-time $O(|V| + |E|)$ algorithm is executed via an explicit work stack rather than recursion, eliminating stack-overflow vulnerabilities on deep structures. Internal self-loops within contracted SCCs are explicitly stripped, guaranteeing a well-founded condensation DAG.

\subsection{Terminal SCCs and Trajectory Confluence}
In dynamic rewriting systems and trajectory execution, deciding whether divergent computational paths eventually converge to a unique terminal attractor is essential.

\begin{definition}[Terminal SCC]
An SCC $C \in V_{\SCC}$ is \emph{terminal} if it has no outgoing edges in $G_{\SCC}$:
\[
\forall u \in C,\quad \forall v \in V,\quad (u, v) \in E \implies v \in C.
\]
\end{definition}

\begin{theorem}[Confluence Criterion via Terminal SCCs]
A finite state transition system represented as an APG $G$ is confluent (every state can reach a common attractor) if and only if its condensation $G_{\SCC}$ contains exactly one terminal SCC.
\end{theorem}
\begin{proof}
If $G_{\SCC}$ has multiple terminal SCCs $C_1 \ne C_2$, then any state entering $C_1$ cannot reach $C_2$ (by definition of terminality), and no path exists from $C_2$ to $C_1$. Thus states in $C_1$ and $C_2$ have no common descendant, refuting confluence. Conversely, if there is exactly one terminal SCC $C^*$, then because $G_{\SCC}$ is a finite DAG, every path from the root must eventually terminate in $C^*$. Hence all trajectories converge into $C^*$, establishing confluence.
\end{proof}

This provides a deterministic decision procedure (\texttt{terminal\_sccs} in \texttt{meta\_calculus.py}) running in linear time.

\section{The Real Analysis Bridge}
While hypersets are discrete graph-theoretical objects, their non-well-founded nature allows them to encode infinite streams, continuous paths, and real numbers via terminal coalgebraic sequences.

\begin{definition}[Hyperset Real Representation]
A real number $x \in [0, 1]$ is represented as an infinite dyadic stream encoded by a cyclic or non-well-founded APG. Rational numbers correspond to finite cyclic graphs (periodic dyadic expansions), while irrational real numbers correspond to infinite paths or computable coalgebraic generators.
\end{definition}

\begin{theorem}[Metric Completeness and Coinductive Continuity]
Let $\mathcal{H}$ be the space of strongly extensional APGs equipped with the bisimulation metric:
\[
d(G_1, G_2) = 2^{-\sup \{k \in \N : G_1 \approx_k G_2\}},
\]
where $\approx_k$ denotes $k$-step bounded bisimulation. Then $(\mathcal{H}, d)$ is an ultrametric complete space, and the projection onto the unit interval $\pi\colon \mathcal{H} \to [0, 1]$ is uniformly continuous.
\end{theorem}
\begin{proof}
Bounded bisimulation $\approx_k$ forms a filtration with $\approx_{k+1} \subseteq \approx_k$. The ultrametric inequality $d(x, z) \le \max(d(x, y), d(y, z))$ follows immediately from the transitivity of $\approx_k$. Cauchy sequences in $\mathcal{H}$ stabilize at each finite depth $k$, defining a unique limit APG under AFA.
\end{proof}

\section{Hypercomputational Oracle Stratification and Verifier Standard Grounding}
How does Grounded Hyperset Theory support the oracle requirements of superintelligences?
\begin{enumerate}
\item \textbf{Autonomous Reflexive Loops:} An agent modeling its own action-perception loop $S_{t+1} = F(S_t, A_t)$ can represent the infinite feedback trajectory as a finite APG with a cyclic attractor, avoiding non-terminating meta-regress.
\item \textbf{Coalgebraic Oracle Simulation:} Higher-order oracles (such as hyperarithmetic truth predicates $\Delta^1_1$) can be characterized coinductively as fixed points of monotone operators over hyperset streams. This is a statement about what the coalgebraic setting can express, not a decision procedure: nothing in the reference implementation decides $\Delta^1_1$ membership.
\item \textbf{Bound-Preserving Evidence, as an intended interface:} Embedding each hyperset operation into the Verifier Standard (VSTD) is the route by which a claim about an agent's self-stability or convergence to an attractor would be carried as an explicit, checkable receipt. In the VSTD-Graph-1..5 assurance architecture, forward support (\texttt{TRUST}), backward diagnostic traversal (\texttt{RUST}), and time degradation (\texttt{ROT}) are defined over graph nodes, and an APG with its SCC condensation is the graph they would be defined over here. \textbf{No such receipt is emitted.} The word ``receipt'' does not occur anywhere in this library's source; there is no receipt machinery, no canonical digest, and no claim record. This item states an interface this framework is built to accept, and no result in this paper rests on it.
\end{enumerate}

\section{Evidential Boundary}
This section states what the preceding sections do and do not establish, so that a reader does not have to infer it.

\noindent\textbf{What is proved.} Four theorems, each with a proof given here: constructive quotient and canonical minimization under Paige--Tarjan refinement, acyclicity of the SCC condensation, the confluence criterion via terminal SCCs, and metric completeness with coinductive continuity. These are prose proofs, and they rest on established results --- Aczel's Anti-Foundation Axiom \cite{aczel}, Paige--Tarjan partition refinement \cite{paige_tarjan}, and Tarjan's SCC algorithm \cite{tarjan}. The mathematics underneath them is not new here. What this paper contributes is the arrangement: the dual-abstraction paradigm that distinguishes bisimulation minimization from SCC condensation, and the bridge from hyperset quotient spaces to computable real analysis.

\noindent\textbf{What is not proved.} There is no proof assistant in this work. This repository carries no Lean, Coq, or Agda development, and no theorem above has been machine-checked. The reference implementation is covered by a test suite that passes, and a passing test suite establishes that the implementation agrees with the specification on the cases tested; it does not establish the specification, and it is not a proof of any theorem here. No claim in this paper should be read as carrying proof-assistant evidence, because none does.

\noindent\textbf{What is out of scope.} Adequacy is not addressed: nothing here shows that APGs with bisimulation quotients are the right representation for self-modeling agents, only that they are a consistent and computable one. The oracle stratification of Section 7 is an account of what the framework can express, not a construction of any hypercomputational object, and no oracle above the arithmetic hierarchy is realized anywhere in the implementation. Receipt emission is future work, and the discipline it would impose --- naming the scope, the limitation and the falsification condition of every result --- is the discipline this section is written in.

\section{Conclusion}
Grounded Hyperset Theory gives a constructive, computable formulation of non-well-founded set theory: APG representations under AFA, polynomial-time bisimulation quotients, linear-time SCC condensation, and a metric completion bridging the quotient space to computable real analysis. Circular self-models and the representational demands of oracle stratification are shown to be treatable in a paradox-free setting, with the standard of evidence stated in Section 8 rather than assumed --- four proofs resting on established algorithmic results, a tested reference implementation, and no machine-checked theorem. The framework is consistent and computable. Whether it is the right framework for self-modeling agents is the open question, and this paper does not settle it.

\begin{thebibliography}{9}
\small
\bibitem{aczel} Peter Aczel. \emph{Non-Well-Founded Sets}. CSLI Lecture Notes, Vol. 14, Stanford University, 1988.
\bibitem{paige_tarjan} Robert Paige and Robert E. Tarjan. Three Partition Refinement Algorithms. \emph{SIAM Journal on Computing} 16(6), 973--989, 1987. \href{https://doi.org/10.1137/0216062}{doi:10.1137/0216062}.
\bibitem{tarjan} Robert E. Tarjan. Depth-First Search and Linear Graph Algorithms. \emph{SIAM Journal on Computing} 1(2), 146--160, 1972. \href{https://doi.org/10.1137/0201010}{doi:10.1137/0201010}.
\bibitem{barwise_moss} Jon Barwise and Lawrence Moss. \emph{Vicious Circles: On the Mathematics of Non-Wellfounded Phenomena}. CSLI Publications, Stanford, 1996.
\bibitem{rutten} Jan J. M. M. Rutten. Universal Coalgebra: A Theory of Systems. \emph{Theoretical Computer Science} 249(1), 3--80, 2000. \href{https://doi.org/10.1016/S0304-3975(00)00056-6}{doi:10.1016/S0304-3975(00)00056-6}.
\bibitem{roost_hypermath} Tyler Roost. \emph{Constructive Quine Fixed-Point Synthesis, Quadrilateral Filtration, and APG Bisimulation}. Preprint, version 0.2.0, 2026.
\end{thebibliography}

\end{document}
